\documentclass{situationist-preprint}

\title{Indications Toward the Unapparent}
\msid{PHE-15-a}
\msfamily{PHE-15}
\msdiscipline{Phenomenology}
\mstype{Principal manuscript}
\mskeywords{inapparent; formal indication; Heidegger; Marion; Henry; Janicaud; method}

\begin{document}
\maketitle

\begin{abstract}
In the seminar at Zähringen in 1973 Heidegger called the thinking he was attempting a phenomenology of the inapparent. The phrase looks self-defeating: phenomenology describes what shows itself, and what does not appear cannot show itself. Janicaud argued that the French phenomenologists who took up the programme, Marion and Henry above all, had in consequence left phenomenology for theology. We argue that a phenomenology of the unapparent is possible, but that its results have a distinctive status: they cannot be shown, only indicated. We distinguish showing from indicating, prove that what is unapparent can be indicated and cannot be shown, and argue that indication is a phenomenological procedure because the indicating relation is itself apparent. The conditions of appearing, the world as horizon and life as self-affection are the principal cases. The result vindicates the programme while limiting what it can claim.
\end{abstract}

\section{The problem}

Phenomenology, in Heidegger's definition, lets that which shows itself be seen from itself, in the very way in which it shows itself from itself \citep[§7]{Heidegger1927}. Yet the same section already says that phenomenology is needed because what it studies \emph{proximally and for the most part} does not show itself: the being of beings is hidden in what does show itself, as its meaning and ground. Late in his career Heidegger radicalized the point. In the Zähringen seminar of 1973 he described the thinking that asks after the presencing of what is present as a phenomenology of the inapparent \citep{Heidegger2003}.

The programme was taken up in France. \citet{Henry1963} argued that the essence of manifestation, life's self-affection, never appears in the horizon of the world. \citet{Marion1997} argued that givenness, which conditions every appearing, is not itself a phenomenon among others. \citet{Janicaud1991} replied that these were not phenomenological results but theological commitments: a method defined by what appears cannot vouch for what does not.

\section{Showing and indicating}

\begin{definition}\label{d:show}
A phenomenological result \emph{shows} $x$ if it describes $x$ as $x$ appears in an intuition that the reader can reactivate.
\end{definition}

\begin{definition}\label{d:indicate}
A phenomenological result \emph{indicates} $x$ if it describes an apparent structure $y$ and an apparent relation by which $y$ refers beyond itself to $x$, without describing $x$ as it appears.
\end{definition}

Indication in this sense is close to the \enquote{formal indication} of Heidegger's early lectures, which points the reader towards a phenomenon to be enacted without fixing it in a content \citep{Heidegger1995}, and to Husserl's horizonal reference, by which every appearance refers to what is co-intended and not given. It differs from Husserl's indicative sign in the \emph{Logical Investigations}, which relates two items both of which can appear.

\begin{definition}\label{d:unapparent}
Something is \emph{unapparent} if it cannot appear as itself.
\end{definition}

\begin{assumption}\label{a:conditions}
Some features of what appears refer beyond themselves to conditions of appearing, and this reference is itself apparent.
\end{assumption}

Assumption~\ref{a:conditions} is supported by three analyses. Every appearance is given within a horizon, and the horizon refers to further horizons and finally to the world, which is the horizon of all horizons and therefore not itself given within one (PHE-11). Every appearance is given to someone as something (PHE-02), and the givenness of each appearance refers to the giving, which is not one more given. And every reduction leaves an act it cannot bracket (PHE-20), whose performing is lived through but not displayed as an object.

\begin{theorem}\label{t:main}
A phenomenology of the unapparent is possible, and its results indicate the unapparent but do not show it.
\end{theorem}

\begin{proof}
By Assumption~\ref{a:conditions}, there are apparent structures with an apparent reference beyond themselves to conditions of appearing, and some of these conditions are unapparent. Describing such a structure and its reference is, by Definition~\ref{d:indicate}, indicating the unapparent; since the description concerns only what appears, it is phenomenological. Showing the unapparent would require describing it as it appears (Definition~\ref{d:show}), and by Definition~\ref{d:unapparent} it does not appear as itself. So the results indicate and do not show.
\end{proof}

\section{Objections}

\begin{objection}[Janicaud]
Indication is a backdoor for theology. Anything can be said to be indicated by something that appears.
\end{objection}

\begin{reply}
Definition~\ref{d:indicate} requires that the indicating relation be apparent. The world's being the horizon of horizons, or the act's not being bracketed by itself, can be checked by anyone who reactivates the corresponding analysis; a transcendent cause of appearing cannot. The definition thus separates the phenomenology of the unapparent from theology on the ground Janicaud himself chose, the accountability of results to what appears.
\end{reply}

\begin{objection}[Henry]
Life's self-affection is not indicated but given, with more certainty than any appearance in the world.
\end{objection}

\begin{reply}
We grant that it is lived. But a result \emph{about} it is a description, and the description presents only the structures of worldly appearing that self-affection is said to underlie. Henry's own analyses proceed by showing what the horizon of the world cannot contain; they indicate life by its non-appearance there.
\end{reply}

\begin{objection}[Marion]
Saturated phenomena show themselves in excess of every concept. The unapparent is thus a matter of excess, not absence.
\end{objection}

\begin{reply}
The saturated phenomenon appears, and so is not unapparent in the sense of Definition~\ref{d:unapparent}. What does not appear on Marion's account is givenness as such, and of it his results speak by indication: from the folds of the given to the giving \citep{Marion1997}. Our thesis therefore applies to Marion's programme where it reaches the unapparent.
\end{reply}

\section{Conclusion}

A phenomenology of the unapparent is possible. It describes apparent structures that refer, apparently, beyond themselves to conditions of appearing that do not appear. Its results are genuine and checkable, but their mode is indication: they point to the unapparent and cannot show it.

\begin{thebibliography}{9}\small
\bibitem[Heidegger(1927)]{Heidegger1927} Heidegger, M. (1927). \emph{Sein und Zeit}. Niemeyer.
\bibitem[Heidegger(1995)]{Heidegger1995} Heidegger, M. (1995). \emph{Phänomenologie des religiösen Lebens} (Gesamtausgabe 60). Klostermann. (Lectures of 1920--21.)
\bibitem[Heidegger(2003)]{Heidegger2003} Heidegger, M. (2003). \emph{Four Seminars} (A. Mitchell \& F. Raffoul, Trans.). Indiana University Press.
\bibitem[Henry(1963)]{Henry1963} Henry, M. (1963). \emph{L'essence de la manifestation}. Presses Universitaires de France.
\bibitem[Janicaud(1991)]{Janicaud1991} Janicaud, D. (1991). \emph{Le tournant théologique de la phénoménologie française}. Éditions de l'Éclat.
\bibitem[Marion(1997)]{Marion1997} Marion, J.-L. (1997). \emph{Étant donné: essai d'une phénoménologie de la donation}. Presses Universitaires de France.
\end{thebibliography}

\end{document}
